\title{LongRoPE: Extending LLM Context Window Beyond 2 Million Tokens}

\begin{abstract}
Expanding the context window in large language models (LLMs) is highly sought after. Nonetheless, significant challenges—including the substantial expenses of fine-tuning, the limited availability of lengthy textual data, and destabilizing effects from previously unseen token positions—constrain current models to context windows of approximately 128k tokens.
\end{abstract}

In this work, we present LongRoPE, a novel approach that, for the first time, enables pre-trained LLMs to handle context windows as large as \textbf{2048k} tokens. This is accomplished using a maximum of 1k fine-tuning steps at sequence lengths up to 256k, all while preserving model performance at the original, shorter context lengths. Our method is built upon three main innovations: (i) we systematically detect and utilize two types of non-uniformities in positional interpolation via an efficient search strategy, which yields superior fine-tuning initialization and permits an 8$\times$ increase in window size without additional fine-tuning; (ii) we propose a stepwise extension mechanism, initially fine-tuning at a 256k context window, followed by a second stage of positional interpolation on the already extended model to attain the 2048k context window; (iii) we further adapt LongRoPE at 8k sequence length to fully restore short context window accuracy. Comprehensive experiments on both LLaMA2 and Mistral across multiple benchmarks confirm the robustness and efficacy of our approach. Notably, LongRoPE maintains the base model architecture, requiring only minimal adjustments to the positional embedding, and remains compatible with most existing system optimizations. Code for LongRoPE will be released at \texttt{https://github.com/microsoft/LongRoPE}

\section{Introduction}

Although Large Language Models (LLMs) have achieved notable advancements across multiple tasks~\cite{gpt4,llama2}, they remain constrained by their finite context window sizes, such as the 4096-token boundary of LLaMA2~\cite{llama2}. When input length surpasses this threshold, model performance deteriorates, primarily because LLMs have not been exposed to those additional token positions during training. This limitation complicates applications that require extensive in-context learning with many examples~\cite{Huang2023FewerIM} and the operation of LLM agents~\cite{park2023generative,madaan2023selfrefine}.

Recent studies demonstrate that the context window of a pre-trained LLM can be expanded up to approximately 128k tokens through fine-tuning on extended texts~\cite{longlora,pi,yarn,zhang2024soaring,liu2023scaling}. However, extending the context window further faces three primary challenges. First, introducing untrained position indices results in a multitude of problematic values, creating out-of-distribution complications and hindering convergence during fine-tuning~\cite{pi}. This issue is exacerbated in cases where the context length is increased from 4k to over 1000k, with more than 90\% of token positions being newly introduced. Second, effective fine-tuning typically depends on the availability of long texts matching the desired context lengths. At present, such ultra-long texts—especially those surpassing 1000k tokens—are scarce within existing datasets. Additionally, training with extremely lengthy sequences is computationally demanding, consuming extensive time and GPU resources. Third, when pushing the context window to extreme lengths, the attention mechanism becomes diluted, distributing focus across an excessive number of tokens and thereby diminishing model performance on original, shorter context tasks~\cite{pi}.

A common strategy to address the initial challenge involves interpolating RoPE positional embeddings~\cite{rope,pi}, where newly encountered position indices are rescaled to align with the pretrained positional range, as illustrated in Fig.~\ref{fig:overview}. Position Interpolation (PI)~\cite{pi} achieves this by performing a linear interpolation of RoPE's rotary angles according to the extension ratio. In contrast, NTK~\cite{ntk,dynamicntk} proposes non-uniform interpolation and extrapolation across the RoPE dimensions. YaRN~\cite{yarn} further divides the RoPE dimensions into three groups based on frequency, applying extrapolation, NTK-based, or linear interpolation strategies to each. Despite these innovations, the positional embeddings within Transformer architectures are characterized by \emph{complex and heterogeneous information entropy}. Existing techniques do not sufficiently capitalize on this nuanced non-uniformity, resulting in information loss and, consequently, constraining the maximum attainable context window.

Section~\ref{sec:analysis} presents two main empirical observations: \textbf{(1)} Successful positional interpolation requires accounting for two types of non-uniformity—specifically, the variation in RoPE dimensions and the distribution of token positions. Minimal interpolation is most beneficial for lower RoPE dimensions and tokens near the sequence start, though the precise optimal strategy is contingent upon the desired sequence extension length. \textbf{(2)} Incorporating awareness of these non-uniformities into positional interpolation helps preserve the essential information encoded in the original RoPE, with particular retention in critical dimensions and positions. This approach mitigates information loss due to interpolation and offers improved initial conditions for subsequent fine-tuning. Additionally, in settings where fine-tuning is not performed, this method supports up to an 8$\times$ increase in sequence length.

In response to these observations, we propose \textbf{LongRoPE}, a robust approach designed to enlarge the LLM context window to accommodate more than 2 \emph{million} tokens. LongRoPE incorporates three principal advancements. First, {LongRoPE} harnesses the full spectrum of multidimensional non-uniformities within positional interpolation. It determines suitable rescale factors for RoPE's rotation angles across each RoPE dimension, contingent on token locations. Given that the search space for these rescale factors increases exponentially with the desired extension ratio, {LongRoPE} utilizes an evolutionary search algorithm enhanced by two optimization strategies to improve search performance. Figure~\ref{fig:overview} provides an illustration of a resulting rescaled RoPE identified through this process.

Subsequently, {LongRoPE} adopts a resource-efficient, stepwise extension method to obtain a 2048k context window, circumventing the necessity for direct fine-tuning on exceptionally long texts, which are both rare and difficult to source. The process initiates with identifying a suitable 256k context length for the pre-trained LLM and performing fine-tuning at this length. Leveraging our non-uniform positional interpolation, which facilitates an 8$\times$ increase in length even without fine-tuning, a subsequent search is carried out to determine new RoPE rescale factors for the already fine-tuned and extended LLM. This approach successfully realizes a 2048k context window for models such as LLaMA2 and Mistral~\cite{mistral}.

In order to prevent reduced performance on the initial (short) context window, {LongRoPE} persists in optimizing the RoPE rescale factors even after expanding the LLM's context. Analogous to the process of scaling up from 256k to 2048k, we apply our search algorithm to decrease the context window to 4k and 8k on the 256k fine-tuned LLM, aiming to reduce excessive positional interpolation. At inference time, when the sequence length falls below 8k, RoPE is modified using the corresponding rescale factors found through this search.

A wide range of experiments conducted on multiple LLMs and diverse long-context tasks confirm the efficacy of our approach. Our results indicate that LongRoPE consistently preserves low perplexity across evaluation lengths spanning from 4k to 2048k, attains passkey retrieval accuracy exceeding 90\%, and offers performance on par with established benchmarks constructed within a 4096-token context window. Furthermore, LongRoPE is compatible with all LLMs that utilize RoPE embedding. We intend to make our codebase and LongRoPE-2048k models publicly available.

\section{Non-uniformity in Positional Interpolation}
\label{sec:analysis}

\subsection{Preliminary}

Transformer architectures depend on clearly defined positional data—commonly realized as position embeddings—to encode the sequence of input tokens. This study concentrates on the RoPE~\cite{rope} position embedding technique, a method that has gained substantial traction in state-of-the-art LLMs. For a token located at position $n$, its RoPE encoding can be succinctly described as:

\begin{equation}
		\fontsize{7.8}{7.8} \selectfont
	\label{eq:rope}
	\left[ cos(n\theta_0), sin(n\theta_0), cos(n\theta_1), \cdot\cdot\cdot, cos(n\theta_{d/2-1}), sin(n\theta_{d/2-1}) \right]   
\end{equation}

Here, $d$ denotes the dimension of the embedding space, 
$n\theta_i$ indicates the rotational angle assigned to the token at position $n$, and 
$\theta_i=\theta^{-2i/d}$ specifies the rotational frequency. Typically, in RoPE, the base value for $\theta$ is set to 10000.

\textbf{\emph{Context window extension ratio $s$} and positional interpolation}. The variable $s$ is introduced as the quotient of the new context window length $L'$ by the initial length $L$, formalized as $s=\frac{L'}{L}$.

In order to expand the context window from $L$ to $L'$, existing positional interpolation techniques propose adjusting the rotation frequencies $\theta_i$ by scaling them down according to the extension ratio $s$. Defining $\beta=\theta^{2/d}$ and using $\lambda$ to represent the specific rescaling parameter associated with $s$, we provide a unified formulation for these positional interpolation approaches as below:

\begin{equation}	
	\fontsize{7.5}{7.5} \selectfont
	\label{eq:unified_rope}
	\begin{aligned}
		\left[cos\left(\frac{n}{\lambda(\beta)^0}\right), sin \left(\frac{n}{\lambda(\beta)^0}\right), cos\left(\frac{n}{\lambda(\beta)^1}\right), \cdot\cdot\cdot,  sin \left(\frac{n}{\lambda(\beta)^{d/2-1}}\right) \right]   
	\end{aligned}
\end{equation}

\textbf{Linear positional interpolation (PI)}. As described in~\cite{pi}, PI employs a straightforward linear interpolation of positional indices when remaining within the bounds of the pre-training sequence length. Given an extension ratio $s$, each positional rotation angle is uniformly scaled down by a factor of $\lambda = s$ in every RoPE dimension. Nonetheless, this results in the compression of position representations, thereby reducing the model's capacity to differentiate between tokens that are near each other in sequence. Consequently, PI often exhibits suboptimal performance when the extension ratio is large.

\textbf{NTK-based interpolation and extrapolation}. The works of~\cite{ntk,dynamicntk} reinterpret RoPE through the lens of information encoding and employ Neural Tangent Kernel (NTK) theory~\cite{jacot2018neural,tancik2020fourier} to analyze its properties. To address the positional clustering seen in PI, these studies propose distributing the task of interpolation more evenly across RoPE’s vector dimensions. Specifically, they reduce the scaling of the lower (higher-frequency) dimensions and amplify scaling in the higher (lower-frequency) ones. This adjustment allows for both interpolation and extrapolation over positions, with the scaling factor $\lambda = s^{i}$. Furthermore, the refined dynamic NTK model~\cite{dynamicntk} adapts the extension rate according to the present sequence length at each position. In contrast to PI, which relies on model fine-tuning, NTK-aware approaches permit context window extension without further fine-tuning, though typically only up to a 4$\times$ increase.

\textbf{YaRN}~\cite{yarn} presents a notable advancement in enhancing the effectiveness of positional interpolation. The approach segments RoPE dimensions into three distinct categories according to their frequency characteristics, with each group receiving a tailored interpolation technique. Specifically, extrapolation ($\lambda$=1) is applied to high-frequency components, while linear interpolation (PI) is designated for those with low frequencies. The remaining intermediate-frequency RoPE dimensions utilize the NTK method. A central feature of YaRN is this frequency-based grouping strategy; however, its current implementation is guided by human intuition and empirical analysis, potentially limiting optimal adaptation to newly developed LLM architectures.

\subsection{Study on Non-uniform Positional Interpolation}

Drawing on insights from NTK and YaRN, we observe that their improvements are primarily due to non-linear treatments, notably through attending to varying frequencies within RoPE dimensions for tailored interpolation and extrapolation. Despite this, existing non-linear approaches are largely based on heuristics crafted by humans. This observation leads to two pertinent questions: (1) Does the prevailing method for positional interpolation achieve optimality? (2) Might there be alternative non-linear strategies that have yet to be investigated?

To investigate these issues, we employ evolution search (refer to Sec.~\ref{sec:method}) to identify improved non-uniform positional interpolations tailored for LLaMA2-7B. Our search process utilizes perplexity as a guiding metric and evaluates candidates using five randomly selected samples from the PG19~\cite{pg19} validation set. Based on our experiments, we present several principal observations.

\textbf{Finding 1}: \textit{RoPE dimensions exhibit substantial non-uniformities, which are not effectively handled by current positional interpolation methods.}

We conduct a search for the optimal value of $\lambda$ corresponding to each RoPE dimension as specified in Eq.~\ref{eq:unified_rope}. Table~\ref{tbl:exp1} presents a perplexity comparison for LLaMA2-7B across different methods on the PG19 and Proof-pile~\cite{proof-pile} test sets, evaluated without any fine-tuning. The results derived from our optimized approach yield notable improvements, indicating that the current strategies—both linear (PI) and non-uniform (Dynamic-NTK and YaRN) interpolation methods—are not optimal. In particular, YaRN exhibits inferior performance compared to PI and NTK on the PG19 dataset, failing to achieve the intended context window length for LLMs not subjected to fine-tuning. As an illustration, when using an 8k context length, YaRN's perplexity escalates sharply beyond 7k tokens.

In our investigation, we observe that the rescaling coefficients $\lambda$ in Eq.~\ref{eq:unified_rope} are not uniform, contrasting with the constant scaling factor $s$ applied in PI, NTK’s formula, and the grouped adjustment used in YaRN. Adopting these non-uniform rescaling factors leads to marked improvements in the language modeling ability of LLaMA2—specifically, reductions in perplexity—when processing context windows of 8k and 16k tokens, all achieved without additional fine-tuning. This enhancement arises because the modified positional embeddings retain the key properties of the original RoPE, especially within critical dimensions, thereby lowering the model's challenge in distinguishing among closely positioned tokens.

\textbf{Finding 2}: \textit{RoPE for the initial tokens in the input sequence should be extrapolated with less interpolation.} 

We propose that the initial $\hat{n}$ tokens within input sequences should undergo reduced interpolation by RoPE, given that these tokens tend to attract higher attention weights—thus playing a pivotal role in the attention mechanism. Such significance is corroborated by prior analyses, notably in Streaming LLM~\cite{streamingllm} and LM-Infinite~\cite{han2023lminfinite}. To empirically investigate this hypothesis, we increase the context length to 8k and 16k using PI and NTK, while selectively disabling interpolation for the first $\hat n$ tokens (evaluated at 0, 2, ..., 256). Notably, with $\hat n=0$, the process defaults to standard PI and NTK. Results presented in Table~\ref{tbl:tokenposition} underscore two principal findings: \textbf{(1)} withholding position interpolation for the initial tokens yields consistent performance improvements; and \textbf{(2)} the optimal choice for $\hat n$ is contingent upon the extent of context window expansion.

\textbf{Finding 3}: \textit{Non-uniform positional interpolation effectively extends LLM context window in both fine-tuning and non-fine-tuning settings.} 

Although our experiments demonstrate that the use of searched non-uniform position interpolation greatly enhances extension capability at 8k and 16k context lengths even without fine-tuning, extending to even longer contexts necessitates fine-tuning. To this end, we fine-tune LLaMA2-7B with our optimized RoPE for a 64k context length (configuration details are provided in the Appendix). As indicated in Table~\ref{tbl:finetune}, our approach achieves superior results compared to PI and YaRN, both prior to and following the fine-tuning of LLaMA2-7B. This improvement stems from leveraging non-uniform positional interpolation, which reduces information degradation and establishes a more effective initialization for fine-tuning.

\textbf{Summary}. Our findings reveal two forms of non-uniformity: differences across RoPE dimensions and across token positions. By capitalizing on these non-uniformities in positional interpolation, we can substantially enhance the performance of LLM context length extension.

\section{LongRoPE}

\label{sec:method}
Building upon these insights, we propose LongRoPE, which initially incorporates an effective search algorithm designed to harness the two identified non-uniformities. Leveraging this approach, LongRoPE is able to extend the context window of LLMs to exceed 2 million tokens.

\subsection{Problem Formulation}

These two forms of non-uniformity significantly expand the solution space and complicate the optimization process. To tackle these challenges, we recast the multidimensional non-uniform position interpolation optimization as a search-based problem.

Considering a large language model designed for a context window of size $L'$, processing long documents $\mathbf{X}$ where every sequence $\mathbf{x} \in \mathbf{X}$ exceeds $L'$ tokens in length, we define the standard rotary angle at the $i^{th}$ dimension for token position $n$ in the RoPE embedding as $\frac{n}{\beta^i}$. Based on this, the optimization objective can be formally expressed as:

\begin{equation}
\begin{gathered}
		\label{eq:problem}

\underset {\mathbf{x}\in\mathbf{X}; \, |\mathbf{x}|\ge L'}{ \operatorname {arg\,min} } \,  \mathcal{L}\, \left( \text{LLM} (\text{RoPE}, \mathbf{X})\right),	\text{where} 
\\
\scalebox{0.93}{
$\underset {\begin{subarray}{l} i=0,\ldots,\frac{d}{2}-1;\\ n\in[0,|\mathbf{x}|);\end{subarray}}{\text{RoPE}_{(n)}}= {\left[\ldots, \cos\left(\mathbb{I}(\hat{\lambda}_{i}, \hat{n}) \cdot \frac{n}{\beta^i}\right), \sin\left(\mathbb{I}(\hat{\lambda}_{i}, \hat{n}) \cdot \frac{n}{\beta^i}\right), \ldots\right] }$}
\\
\text{with} \quad \mathbb{I}(\hat{\lambda}_{i},\hat{n}) = \begin{cases}
1 &\text{if}\;  n < \hat{n} \\
{\frac{1}{\lambda}_{i}} &\text{if}\; n \geq \hat{n}
\end{cases}
\end{gathered}
\end{equation}

In this context, we define a set of rescaling coefficients, $\mathbb{I}(\hat{\lambda}_{i},\hat{n})$, to account for the two distinct types of non-uniformities. Here, $\hat{\lambda}_i$ corresponds to the non-uniformity present across the RoPE dimensions, while $\hat{n}$ refers to that associated with token positions. In particular, the rescaling function $\mathbb{I}(\hat{\lambda}_{i},\hat{n})$ is employed to modify the rotation angle in the $i^{th}$ RoPE dimension, applying the factor $\hat{\lambda}_i$ only after a specific positional threshold $\hat{n}$ is reached. For the first $\hat{n}$ token positions, the original RoPE rotary angle $\frac{n}{\beta^i}$ is maintained without adjustment from $\hat{\lambda}_i$. When the token position $n$ meets or exceeds $\hat{n}$, the rescaling factor is incorporated.

With a desired context window length of $L'$, the aim is to determine the best rescale factors ($\mathbb{I}(\hat{\lambda}_{0},\hat{n})$, $\mathbb{I}(\hat{\lambda}_{1},\hat{n})$, ... $\mathbb{I}(\hat{\lambda}_{i},\hat{n})$, ...) corresponding to each RoPE dimension from the $1^{st}$ to the $d^{th}$. This optimization allows the modified $\text{LLM}$, incorporating the adjusted $\text{RoPE}$, to minimize the next token prediction loss $\mathcal{L}$ (i.e., perplexity) on input sequences $\mathbf{X}$ of token length $L'$.

\subsection{Searching the Non-uniform Position Interpolation}

To address the issue formulated in Eq.~\ref{eq:problem}, we propose an intuitive yet powerful approach that identifies optimal RoPE rescale factors, aiming to capture and leverage multidimensional variations within position embeddings.

\textbf{Search space}. Our method constructs an extensive search space to accommodate the two forms of non-uniformity. The structure of this search space is detailed in Table~\ref{tbl:searchspace}. Rather than seeking $\mathbb{I}(\hat{\lambda}_{i},\hat{n})$ directly, we streamline the process by searching over the parameters $\lambda_i$ and $\hat{n}$, where $\hat{\lambda}_{i}=1/\lambda_i$, and assigning a distinct rescale factor to each dimension within the RoPE embedding. As indicated in Table~\ref{tbl:searchspace}, the value of $\lambda_i$ can be explored within the interval from 1.0 (equivalent to direct extrapolation) up to $s\times1.25$ (allowing for greater interpolation than standard PI), incremented by 0.01. Here, $s$ denotes the desired context window extension ratio.

The parameter $\hat{n}$ determines how many leading token positions maintain their original RoPE embedding, without applying position interpolation. In practice, we set $\hat{n}$ to take values from the set \{0, 1, 2, 4, 8, 12, 16, 20, 24, 28, 32, 64, 128, 256\}. Setting $\hat{n}=0$ indicates that every token position utilizes the rescale factors discovered during the search process.

\textbf{Evolution-based search}. As outlined in Table~\ref{tbl:searchspace}, our search space encompasses a vast array of positional interpolation strategies, making thorough exploration computationally intensive. For instance, when $s=4\times$, the number of potential configurations is $400^{128/2}\times14 = 4\times10^{167}$, and this number grows exponentially with increasing extension ratios. To efficiently navigate this immense space, we adopt evolution search~\cite{guo2020single} and implement two specific optimization strategies that substantially enhance the speed of the search. The complete procedure is detailed in Algorithm~\ref{alg:search}.

\textit{Optimized initial population generation}. Rather than relying solely on random initialization for the population of $P$ rescale factors, we explicitly include the three RoPE rescale factors used by PI, NTK, and YaRN as part of the starting population. The other $P$-3 members of the population are created by applying random mutations to these three rescale factors, each mutation occurring with probability $p$.

\textit{Monotonically non-decreasing constraint}. Once the initial population has been created, we evaluate the LLM perplexity for every individual. This involves assigning the associated RoPE rescale factors to the target LLM and then calculating the perplexity over the input $\mathbf{X}$. The k individuals with the lowest perplexity are chosen as parents for the next evolutionary cycle. Nonetheless, the extensive search space can result in unproductive mutations and crossovers, yielding suboptimal candidates and causing excessive perplexity evaluations. This inefficiency is exacerbated when $L'$ is large, due to the high computational cost of each perplexity inference.

To tackle this issue, we enforce a monotonic non-decreasing constraint on the sequence of RoPE rescaling factors such that $\lambda_i \le \lambda_{i+1}$ holds for all $i$. Only those RoPE configurations adhering to this condition are utilized during LLM perplexity assessment, thus substantially lowering the computational overhead involved in searching. More precisely, the rescaling factors $\lambda_i$ are mandated to grow monotonically with respect to the RoPE dimension index ($i=0,...,63$). This requirement is motivated by NTK theory~\cite{jacot2018neural,tancik2020fourier,ntk}, which posits that lower-index (high-frequency) dimensions are best served with minimal interpolation (smaller $\lambda_i$), whereas higher-index (low-frequency) dimensions benefit from increased interpolation (larger $\lambda_i$).

\begin{algorithm}[t]
	\small
	\caption{The search algorithm}
	\textbf{Input:} target LLM, input samples $\mathbf{X}$, population size $P$, mutation size $N_1$, crossover size $N_2$, max iterations $\mathcal{T}$, mutate probability $p$\\

	\begin{algorithmic}[1]
		\label{alg:search}
		\STATE Initialize Top-k as empty set;	
		\STATE $\text{P}_0$ = \textit{Initialize\_population\_with\_optimization}($P$, $\mathbf{X}$, $p$); 
		\FOR{$i$ = $1$ to $\mathcal{T}$}
		\STATE \textit{Compute\_perplexity} (LLM, $\text{P}_{i-1}$, $\mathbf{X}$);
		\STATE  Top-k $\leftarrow$ \textit{Update\_Topk}(Top-k, $\text{P}_{i-1}$);
		\STATE$\text{P}_{mutation}$ = \textit{Mutation\_with\_mono\_constraint}(Top-k, $N_1$, $p$); 
		\STATE $\text{P}_{crossover}$ = \textit{Crossover\_with\_mono\_constraint}(Top-k, $N_2$);
		\STATE $\text{P}_i$ = $\text{P}_{mutation}$ $\cup$ $\text{P}_{crossover}$ $\cup$ Top-k;
		\ENDFOR
		\STATE Output the individual in Top-k with minimal perplexity;
	\end{algorithmic}
\end{algorithm}

\textbf{8$\times$ extension without fine-tuning}. Through evolutionary search, our approach successfully discovers non-uniform RoPE rescaling factors that safeguard crucial dimensions and positional information, thus reducing information loss due to interpolation. As illustrated in Fig.\ref{fig:non-ft}, this strategy enables extending the LLaMA2 context window from 4k up to 32k tokens without the need for fine-tuning. Conversely, prior methods—including PI, and non-uniform NTK and YaRN—experience significant increases in perplexity when exceeding a 2$\times$ context window extension.

\subsection{Extending LLM Context Window to 2048K}
\label{sec:extendto2048k}

\textbf{Progressive extension to 2048k}. In this section, we present our approach for expanding the context length of pre-trained LLMs from the standard 4k up to beyond 2048k. As previously shown, our non-uniform positional interpolation method enables an 8$\times$ increase in window size without the need for any fine-tuning. When attempting much larger expansions (such as 512$\times$), however, fine-tuning becomes essential. A typical strategy involves identifying RoPE rescaling parameters suitable for the target 2048k context length before initiating fine-tuning. Nevertheless, this process is hindered by the extremely high computational costs involved in training. In addition, our practical experience indicates that effective fine-tuning of LLMs at such significant extension ratios is difficult to achieve (refer to the Appendix for further details).

Fortunately, LongRoPE demonstrates robust performance when applied to both base models and those that have undergone fine-tuning with extended context windows. In light of this, we propose a streamlined, incremental strategy that reaches the desired 2048k context length in merely 1,000 fine-tuning iterations, utilizing training sequences capped at 256k tokens.

$\Diamond$ \textit{LongRoPE search for scaling pre-trained LLMs to 256k}. Using LLaMA2 as a case study, we perform a search aimed at achieving context window lengths of 128k and 256k. This results in extension ratios of 32$\times$ and 64$\times$ at this phase.

$\Diamond$ \textit{Fine-tuning to 256k}. Subsequently, we adapt the pre-trained LLM to support a context window of 256k tokens through additional fine-tuning. The procedure begins by fine-tuning LLaMA2 for 400 iterations employing RoPE rescaled factors set for 128k. Upon completing this stage, we switch the RoPE rescaled factors to accommodate 256k on the resulting checkpoint and proceed with 600 further fine-tuning steps. This staged approach yields greater efficiency compared to attempting to fine-tune directly to 256k.

$\Diamond$ \textit{Scaling fine-tuned extended LLM to 2048k using LongRoPE search}. In the last stage, we apply an additional search procedure to the LLM already fine-tuned on 256k sequence length. This process enables us to expand the model’s context window to a remarkable 2048k tokens, accomplished without any extra fine-tuning steps. Consequently, the model achieves a total extension ratio of $512\times$.

\textbf{Degradation in original context window}. Upon expanding the context window to an extensive 2048k tokens, we observe that the model's accuracy within the initial context range diminishes. This phenomenon has been previously documented with positional interpolation~\cite{pi}, where position embeddings representing higher indices inside the baseline context window are compressed into a constrained segment, leading to reduced model effectiveness. Specifically, when applying a 512$\times$ extension, the position indices within the standard 4k context window are densely packed together.

To address this issue, we conduct an additional evolution search on the extended LLM, specifically tuning the RoPE rescale parameters for shorter context lengths such as 4k and 8k. For these shorter sequences, we constrain the upper bound of the $\lambda$ values considered in the search, as less positional interpolation is necessary. At inference time, the LLM adaptively selects appropriate RoPE rescale factors based on the context length.

\section{LongRoPE}

\label{sec:method}
Building on these insights, we propose LongRoPE—a framework that initially develops an optimized search algorithm designed to leverage both non-uniformities to their fullest potential. Subsequently, this approach is employed to expand the LLM context window to over 2 million tokens.

\subsection{Problem Formulation}

The presence of these two types of non-uniformity expands the range of potential solutions and complicates the optimization process. To overcome this challenge, we recast the multidimensional non-uniform position interpolation optimization into a search-based framework.

Given an LLM designed for a context window of size $L'$, and input documents $\mathbf{X}$ such that each sequence $\mathbf{x}\in\mathbf{X}$ exceeds $L'$ tokens, we represent the original rotary angle for the $i^{th}$ dimension in RoPE at the token position $n$ as $\frac{n}{\beta^i}$. Based on this setup, the optimization problem can be described as follows:

\begin{equation}
\begin{gathered}
		\label{eq:problem}

\underset {\mathbf{x}\in\mathbf{X}; \, |\mathbf{x}|\ge L'}{ \operatorname {arg\,min} } \,  \mathcal{L}\, \left( \text{LLM} (\text{RoPE}, \mathbf{X})\right), \quad \text{with} 
\\
\scalebox{0.93}{
$\underset {\begin{subarray}{l} i=0,\cdot\cdot,\frac{d}{2}-1;\\ n\in[ 0,|\mathbf{x}|);\end{subarray}}{\text{RoPE}_(n)}= {\left[\cdots,\cos\left(\mathbb{I}(\hat{\lambda}_{i}, \hat{n})\cdot\frac{n}{\beta^i}\right), \sin\left(\mathbb{I}(\hat{\lambda}_{i}, \hat{n})\cdot\frac{n}{\beta^i}\right),\cdots\right] }$}\\
	\text{with} \quad \mathbb{I}(\hat{\lambda}_{i},\hat{n})=\begin{cases}
	1&\mbox{\quad if $n<\hat{n}$}\\
{\frac{1}{\lambda}_{i}}&\mbox{\quad if $n\geq\hat{n}$}
\end{cases}
	\end{gathered}
\end{equation}

We define a set of rescaling coefficients, $\mathbb{I}(\hat{\lambda}_{i},\hat{n})$, designed to address both types of non-uniformity. Here, $\hat{\lambda_i}$ represents the dimension-specific non-uniformity within RoPE, while $\hat{n}$ signifies the positional threshold among tokens. More concretely, $\mathbb{I}(\hat{\lambda}_{i},\hat{n})$ is employed to adjust the rotation angle in the $i^{th}$ RoPE dimension; in this context, $\hat\lambda_{i}$ serves as the scaling parameter, and $\hat{n}$ designates when this adjustment comes into effect. For the first $\hat{n}-1$ token positions, the scaling factor $\hat\lambda_{i}$ remains inactive and the standard rotary angle $\frac{n}{\beta^i}$ for RoPE is applied. Once token positions reach $n \ge \hat{n}$, the rescaling factor is activated.

Assuming a desired context window of length $L'$, our aim is to determine the best rescaling coefficients ($\mathbb{I}(\hat{\lambda}_{0},\hat{n})$, $\mathbb{I}(\hat{\lambda}_{1},\hat{n})$, ... $\mathbb{I}(\hat{\lambda}_{i},\hat{n})$, ...) corresponding to each RoPE dimension from the $1^{st}$ through the $d^{th}$. This adjustment allows the underlying $\text{LLM}$, equipped with the modified $\text{RoPE}$, to attain the lowest possible next-token prediction loss, denoted as $\mathcal{L}$ (equivalently, minimizing perplexity), on input sequences $\mathbf{X}$ of length $L'$.

\subsection{Searching the Non-uniform Position Interpolation}

To address the problem formulated in Eq.~\ref{eq:problem}, we present a straightforward yet powerful approach that determines the optimal RoPE rescale factors, thereby leveraging the complex multidimensional non-uniformities present in position embedding.

\textbf{Search space}. We construct an extensive search space that captures both types of non-uniformities. An overview of this search space is provided in Table~\ref{tbl:searchspace}. Specifically, our method involves assigning a distinct rescale factor to each dimension within the RoPE embedding. For clarity and simplicity in defining the search space, we perform the search over $\lambda_i$ and $\hat{n}$ rather than directly over $\mathbb{I}(\hat{\lambda}_{i},\hat{n})$, where the relation $\hat{\lambda}_{i}=1/\lambda_i$ holds. As illustrated in Table~\ref{tbl:searchspace}, the value of $\lambda_i$ can range from a minimum of 1.0 (equivalent to standard extrapolation) up to a maximum of $s\times1.25$ (which corresponds to a greater degree of interpolation compared to PI), with increments of 0.01. Here, $s$ denotes the intended ratio for extending the context window.

The parameter $\hat{n}$ specifies how many of the earliest token positions are preserved using the unaltered RoPE embedding, foregoing position interpolation. In practice, $\hat{n}$ is selected from the set \{0, 1, 2, 4, 8, 12, 16, 20, 24, 28, 32, 64, 128, 256\}. If $\hat{n}=0$, this indicates that every token position relies on the optimized rescale factors.

\textbf{Evolution-based search}. The search space outlined in Table~\ref{tbl:searchspace} encompasses a vast array of positional interpolation configurations, making exhaustive exploration computationally infeasible. For instance, with an extension ratio of $s=4\times$, the number of possible configurations is $400^{128/2}\times14$=4$\times10^{167}$. This exponential growth in search space with increasing extension ratios poses additional difficulties. To overcome these challenges, we employ evolution search~\cite{guo2020single} along with two optimization strategies designed to significantly enhance the efficiency of the search process. The complete procedure is presented in Algorithm~\ref{alg:search}.

\textit{Optimized initial population generation}. Rather than relying solely on random selection for the initial set of $P$ rescale factors, the initialization includes the three RoPE rescale factors associated with PI, NTK, and YaRN as predefined members of the population. The remaining $P-3$ individuals are generated by introducing random mutations to these three rescale factors, applying each mutation with probability $p$.

\textit{Monotonically non-decreasing constraint}. Following the initialization of the population, LLM perplexity is evaluated for every candidate. In this process, we apply the designated RoPE rescale factors within the target LLM and measure the perplexity on input $\mathbf{X}$. The $k$ individuals with the lowest perplexity are then selected as parents for subsequent evolutionary steps. Nonetheless, the large size of the search space means that simple mutation and crossover can often traverse suboptimal regions, resulting in superfluous perplexity evaluations. This inefficiency is exacerbated as $L'$ increases, since each perplexity computation demands significant inference time.

To tackle this issue, we enforce a constraint that ensures the RoPE rescaled factors sampled are non-decreasing, namely $\lambda_i \le \lambda_{i+1}$. Only RoPE configurations that meet this monotonicity criterion are utilized for perplexity assessment in the LLM, thereby streamlining the search space considerably. In detail, we stipulate that $\lambda_i$ must exhibit a monotonically increasing trend across the RoPE dimensions (that is, $i$ = 0,...,63). This requirement for dimensional monotonicity is motivated by insights from NTK theory~\cite{jacot2018neural,tancik2020fourier,ntk}, which indicate that lower-indexed dimensions (representing higher frequencies) necessitate less interpolation (i.e., smaller $\lambda_i$ values), whereas higher-indexed dimensions (corresponding to lower frequencies) benefit from greater interpolation (i.e., larger $\lambda_i$ values).

\begin{algorithm}[t]
	\small
	\caption{The search algorithm}
	\textbf{Input:} target LLM, input samples $\mathbf{X}$,   population size $P$, mutation size $N_1$, crossover size $N_2$, max iterations $\mathcal{T}$, mutate probability $p$\\

	\begin{algorithmic}[1]
		\label{alg:search}
		\STATE Top-k $\leftarrow \phi$;	
		\STATE $\text{P}_0$ $\leftarrow$ \textit{Initialize\_population\_with\_optimization} ($P$, $\mathbf{X}$, $p$); 
		\FOR{$i$ from $1$ to $\mathcal{T}$}
		\STATE \textit{Compute\_perplexity} (LLM, $\text{P}_{i-1}$, $\mathbf{X}$);
		\STATE Top-k $\leftarrow$ \textit{Update\_Topk} (Top-k, $\text{P}_{i-1}$);
		\STATE $\text{P}_{mutation}$ $\leftarrow$ \textit{Mutation\_with\_mono\_constraint} (Top-k, $N_1$, $p$); 
		\STATE $\text{P}_{crossover}$ $\leftarrow$ \textit{Crossover\_with\_mono\_constraint} (Top-k, $N_2$);
		\STATE $\text{P}_i$ $\leftarrow$ $\text{P}_{mutation} \cup \text{P}_{crossover} \cup$ Top-k;
		\ENDFOR
		\STATE Output the member of Top-k exhibiting the lowest perplexity;
	\end{algorithmic}
\end{algorithm}

\textbf{8$\times$ extension without fine-tuning}. Leveraging evolutionary search, our approach efficiently discovers optimal non-uniform RoPE rescale parameters, selectively retaining critical dimensions and positional information to reduce information loss due to interpolation. As illustrated in Fig.\ref{fig:non-ft}, this enables us to expand LLaMA2's context window from 4k up to 32k tokens without requiring fine-tuning. By comparison, other approaches—including PI as well as non-uniform NTK and YaRN—result in a significant increase in perplexity beyond a 2$\times$ extension.

\subsection{Extending LLM Context Window to 2048K}
\label{sec:extendto2048k}

\textbf{Progressive extension to 2048k}. In this section, we present our approach for scaling the context window of pre-trained LLMs from the standard 4k up to more than 2048k tokens. As previously shown, our non-uniform positional interpolation technique enables an extension by a factor of 8 without the need for additional fine-tuning. However, achieving even greater expansions (such as a 512$\times$ increase) necessitates fine-tuning. One possible strategy involves identifying suitable RoPE rescaling parameters for the target 2048k window size, followed by fine-tuning. Nevertheless, this approach is hindered by the substantial computational resources required for training. Additionally, our observations indicate that effective fine-tuning becomes increasingly difficult when attempting very large context window extensions (refer to Appendix for further discussion).

Fortunately, LongRoPE demonstrates efficacy when applied to both the base and the fine-tuned extended LLMs. Accordingly, we propose a resource-efficient and incremental approach capable of reaching the desired 2048k target using only 1k fine-tuning steps and a maximum training length of 256k.

$\Diamond$ \textit{LongRoPE search for adapting pre-trained LLMs to 256k}. Using LLaMA2 as a case study, we perform a search procedure aimed at achieving context window sizes of 128k and 256k. In this phase, the context window is extended by factors of 32$\times$ and 64$\times$, respectively.

$\Diamond$ \textit{Fine-tuning to 256k}. In the next stage, the pre-trained LLM is further fine-tuned to support a 256k context window. The process begins by fine-tuning LLaMA2 for 400 steps utilizing the RoPE rescaled factors set for 128k. After this, we swap the RoPE rescaled factors to those corresponding to 256k in the obtained checkpoint, followed by an extra 600 steps of fine-tuning. This approach demonstrates greater efficiency compared to direct fine-tuning at 256k.

$\Diamond$ \textit{Expanding fine-tuned extended LLMs to 2048k through LongRoPE search}. In the final step, we conduct an additional search using the fine-tuned LLM with a 256k context length. This process yields an exceptionally large context window of 2048k, achieved without any further fine-tuning. As a result, the total extension factor reaches $512\times$.

\textbf{Performance degradation within original context window}. Upon increasing the context window size to 2048k, we observe a decrease in model performance when evaluated on the original window size. This phenomenon, well-documented in prior research on positional interpolation~\cite{pi}, arises because position embeddings corresponding to the original context are now mapped to a significantly compressed region of the embedding space in higher dimensions, which impairs the language model’s effectiveness. In particular, an extension ratio of 512$\times$ results in dense clustering of positions within the initial 4k context window.

To address this issue, an additional evolutionary search is conducted on the extended LLM, specifically tuning the RoPE rescale factors for shorter context lengths such as 4k and 8k tokens. For these shorter contexts, the upper bound for the searched $\lambda$ is decreased, reflecting the reduced need for positional interpolation. At inference time, the LLM adaptively applies the appropriate RoPE rescale factors as needed.

 \section{Experiments}

\label{sec:exp}
\subsection{Setup}
\textbf{Evaluation Tasks and Models}. LongRoPE is integrated with LLaMA2-7B and Mistral-7B, and assessed across three criteria: (1) perplexity on long-form documents to gauge the models' handling of extended contexts; (2) performance on the Passkey retrieval task, designed to test the capability to extract a specific passkey from predominantly distractor text; and (3) results on established LLM benchmarks using the conventional 4096-token context window.

\textbf{Fine-tuning}. For LLaMA2, we set the learning rate at 2e-5, employing a linear decay schedule and a total batch size of 32 across all devices. The model undergoes 400 fine-tuning steps on the Redpajama~\cite{redpajama} dataset, divided into 128k-length segments, each demarcated by the BOS and EOS tokens. Subsequently, starting from the resulting checkpoint, training is continued for an additional 600 steps to extend the context window to 256k tokens. Fine-tuning at the 128k context size utilizes 8 A100 GPUs with the distributed training framework~\cite{cube}, whereas the 256k context length necessitates 16 A100 GPUs. For Mistral, a fixed learning rate of 1e-6 is adopted with an overall batch size of 64. Training for both the 128k and 256k variants mirrors the approach in YaRN~\cite{yarn}, with each undergoing 400 steps on Together Computer's Long-Data Collections~\cite{mistraldata} at a sequence length of 16k. All Mistral experiments are conducted using 4 A100 GPUs.

\textbf{Search}. When the target window size does not exceed 256k, the following parameters are used: $P$=64, $N_1$=$N_2$=16, $p$=0.3, $\mathcal{T}$=40, and the top 32 candidates are chosen for mutation and crossover in every iteration. Perplexity is evaluated on 5 randomly selected samples from the PG19 validation set, each meeting at least the desired context length. For window sizes exceeding 512k, both the population size and the mutation/crossover counts are reduced by half. In this scenario, perplexity assessment is carried out using 3 randomly chosen validation samples from Pile-Books3~\cite{pile}.

\textbf{Baselines}. In order to extend to a 2048k context window, we performed fine-tuning on models with 128k and 256k contexts. This produces LongRoPE-2048k (ft=128k) and LongRoPE-2048k (ft=256k) variants for LLaMA2 and Mistral, respectively. We benchmark these four models against leading baselines for large-context window adaptation, focusing on publicly available LLMs that have undergone positional interpolation fine-tuning through PI, NTK, and YaRN methods. Our baseline set includes Together-32k~\cite{together}, Code LLaMA~\cite{codellama}, LongLoRA-full-FT-100k~\cite{longlora}, as well as YaRN-LLaMA and YaRN-Mistral~\cite{yarn}.

\subsection{Main Results}

\label{sec:mainresult}
\textbf{Long sequence language modeling within 256k}. We start by benchmarking against leading extended LLMs on sequences up to 256k tokens in length. To illustrate our model's generalizability, we employ two datasets: the Proof-pile~\cite{pg19} and PG19~\cite{pile} test sets. Perplexity is measured across a range of context lengths utilizing a sliding window approach with a window size of 256. The full test set of PG19, containing 100 documents, is used for evaluation. For Proof-pile, we adopt the selection protocol from YaRN~\cite{yarn}, randomly choosing 10 samples, each with a minimum length of 128k tokens.

Table~\ref{tbl:proof-pile} and Table~\ref{tbl:pg19} present a comparison of the perplexity scores for LLaMA2 and Mistral models, each extended using various interpolation approaches, on the Proof-pile and PG19 datasets, respectively. Two major findings are noteworthy: \textbf{(1)} the extended models consistently exhibit a reduction in perplexity as the evaluation length increases from 4k to 256k, indicating their enhanced capability to utilize longer contexts. \textbf{(2)} Remarkably, even when evaluated with a context window that is 16$\times$ greater—a condition typically associated with degraded performance at shorter lengths—our LongRoPE-2048k models achieve lower perplexity compared to leading baseline methods for context lengths up to 256k.

\textbf{Long sequence language modeling beyond 2000k}. To assess model performance on exceptionally lengthy texts, the Books3~\cite{pile} dataset is utilized. For efficient evaluation, we randomly sample 20 books, all longer than 2048k tokens, applying a 256k sliding window during testing.

As detailed in Table~\ref{tbl:books3}, LongRoPE is able to expand the context window of both LLaMA2-7B and Mistral-7B models up to 2048k, while still maintaining perplexity at levels comparable to or even better than those of baseline models for shorter input lengths between 8k and 128k. Distinct performance disparities emerge when comparing the 2048k configurations of LLaMA2 and Mistral. Specifically, Mistral shows superior results relative to baselines at shorter contexts; however, its perplexity rises above 7 when the context length surpasses 256k. The LLaMA2 model behaves as anticipated: perplexity generally improves (decreases) as the context grows longer, though there are slight increases noted at the 1024k and 2048k intervals. Further, for LLaMA2, utilizing a fine-tuning window of 256k for LongRoPE-2048k produces better outcomes than using 128k, which is attributable to the lower secondary extension ratio (8$\times$ versus 16$\times$). By contrast, Mistral shows optimal performance with a fine-tuning length of 128k. This can be explained by the fact that, following the YaRN configuration, both 128k and 256k fine-tuning on Mistral employ a training sequence length of 16k, thereby constraining Mistral's capacity to scale its context window beyond fine-tuning.

\textbf{Passkey retrieval}. In this section, we examine how context window size influences performance in generation-based tasks. We adopt the synthetic passkey retrieval evaluation introduced by ~\cite{passkey}, where the objective is for the model to locate a concealed random passkey (specifically, a five-digit numeral) within a lengthy text. The structure of the prompts utilized is provided in the appendix. To assess performance, we conduct the passkey retrieval task over 10 trials, each time embedding the passkey at a randomly chosen position, uniformly distributed throughout the entire context window used in evaluation.

Figure~\ref{fig:passkey} presents a comparison of retrieval accuracy against baseline approaches. The accuracy of prior models plummets to zero once the input exceeds 128k tokens. In stark contrast, our LongRoPE-LLaMA2-2048k (ft=256k) maintains strong retrieval accuracy ($\ge$90\%) across the full span from 4k to 2048k tokens, successfully handling the demanding task of recovering a passkey among millions of tokens. Meanwhile, LongRoPE-Mistral-2048k (ft=128k) preserves perfect accuracy up to 1800k tokens, with a decrease to 60\% at 2048k—a result consistent with the findings in Table~\ref{tbl:books3}, where a slight uptick in perplexity is observed at the 2048k token length.

\textbf{Standard benchmarks within original context window}. The LongRoPE-2048k models are assessed on tasks confined to the initial context window, employing the Hugging Face Open LLM Leaderboard~\cite{huggingfaceboard} for evaluation in both zero-shot and few-shot paradigms. Specifically, the 25-shot ARC-Challenge~\cite{allenai:arc}, 10-shot HellaSwag~\cite{zellers2019hellaswag}, 5-shot MMLU~\cite{mmlu}, and 0-shot TruthfulQA~\cite{lin2021truthfulqa} datasets are utilized.

As presented in Table~\ref{tbl:commonsense}, our models attain results similar to those on the initial benchmark, which was intended for use with a limited context window; notably, they surpass the original Mistral on TruthfulQA by a margin of +0.5\%. While LongRoPE-LLaMA2-2048k, after fine-tuning at 256k, exhibits a minor additional decline in performance, its outcomes for the majority of tasks remain acceptably close to baseline.

\subsection{Ablation Results}

\textbf{Effectiveness of the second positional interpolation}. Within our progressive extension approach, we apply our search algorithm to perform a second non-uniform positional interpolation on LLMs that have already been fine-tuned and extended. To assess the impact of this method, we conduct experiments using our fine-tuned LLaMA2-256k model, which we further extend to 512k, 1024k, and 2048k token contexts utilizing PI and YaRN techniques. As illustrated in Table~\ref{tbl:secondsearch}, our non-uniform positional interpolation consistently maintains stable perplexity levels. Conversely, models extended with PI and YaRN demonstrate a marked increase in perplexity as the extension ratio grows.

\textbf{Effectiveness of recovery at shorter context lengths.} To address diminished performance at shorter contexts, we modify the RoPE scaling parameters in LongRoPE-2048k using our search procedure. In particular, the maximum permitted scale factors in the search are reduced, promoting minimal interpolation for shorter 4k and 8k input sequences. 
Table~\ref{tbl:shortperformance} provides a perplexity comparison of LongRoPE-LLaMA2-2048k on the Proof-pile dataset at 4k and 8k tokens, as well as the mean accuracy across LLM benchmarks. These findings clearly indicate notable gains in performance at reduced context lengths.

\textbf{Analysis on the two forms of non-uniformities}. To further investigate the role of the two types of non-uniformities, we conduct an ablation study to assess the individual impact of each component on model performance. Specifically, we design two experimental setups: (i) scaling LLaMA2-7B to shorter sequence lengths of 16k and 32k by applying several approaches—PI, searching exclusively over the RoPE dimension, and searching over both sources of non-uniformity; (ii) increasing the context length of a previously fine-tuned LLaMA2 model (trained at 256k tokens) to 2048k using the same set of strategies. Perplexity scores are measured without further fine-tuning. Results reported in Table~\ref{tbl:searchdimension} indicate that introducing non-uniformity along the RoPE dimension yields a considerable decrease in perplexity relative to the linear interpolation baseline of PI. Additionally, incorporating non-uniformity in token positions provides notable gains at the 16k and 32k levels, although such benefits do not extend to the 2048k setting, which may be attributable to the challenge posed by the extraordinary sequence length. Simply retaining initial tokens without interpolation appears ineffective under these circumstances; we propose to explore this observation in future research.

\section{Related Works}

Beyond approaches reliant on position interpolation, this section reviews alternative strategies within the literature.

\textbf{Retrieval-based} approaches employ an external memory component to store earlier contextual information, coupled with retrieval mechanisms to access pertinent documents during inference~\cite{longllama,wang2023augmenting,borgeaud2022improving}. These approaches generally require significant changes to the underlying LLM architectures. By comparison, our method remains relatively lightweight, involving only modest adjustments to positional embeddings. Furthermore, our approach supports a wider range of long context applications beyond retrieval, including long-form document summarization and few-shot learning.

\textbf{Attention-based context window extensions}. In addition to positional embedding interpolation, several studies have proposed expanding the input context using the standard LLM context window, primarily through modifications of attention mechanisms~\cite{han2023lminfinite, streamingllm,parallelwindow}. Central to these methods is the introduction of innovative attention masks to address the challenge of attention explosion arising from added positions. These approaches serve as complements to positional interpolation techniques.

\textbf{Fine-tuning based approaches} investigate strategies to adapt pre-trained LLMs for extended context windows by modifying position embeddings. Methods such as Code LLaMA~\cite{codellama}, LLaMA2 Long~\cite{xiong2023effective}, and ScaledRoPE~\cite{liu2023scaling} adopt significantly increased RoPE base values and perform fine-tuning at the intended target sequence lengths. In contrast, our approach accommodates a range of target lengths and is capable of reaching beyond 2M tokens. More recently, the substantial GPU resources required for fine-tuning models at long context lengths (over 128k tokens) have motivated techniques like LongLoRA~\cite{longlora} and PoSE~\cite{zhu2023pose} to reduce such computational demands. Our proposed method remains complementary to these efficient fine-tuning strategies.

\section{Conclusion}

In this study, we introduce LongRoPE, a technique that dramatically increases the context length of LLMs to 2048k, all while preserving their performance in the original, shorter context windows. Our approach leverages two distinct non-uniform aspects in RoPE positional encoding, identified via an efficient evolutionary search algorithm. This method yields two significant advantages: it furnishes strong initialization for subsequent fine-tuning, and it facilitates an 8$\times$ expansion of the context window in the absence of fine-tuning. Building upon these results, we develop a progressive extension approach that employs LLMs fine-tuned on 256k-length contexts to attain a 2048k context window, without requiring additional fine-tuning steps. A comprehensive suite of experiments confirms the efficacy of LongRoPE. We anticipate that the LongRoPE-2048k models will open avenues for novel applications requiring very long context and stimulate new directions in research.

\section*{Broader Impacts} This study aims to contribute to progress in the field of Machine Learning. While our research may carry various societal implications, we do not identify any particular issues that warrant explicit mention in this context.

	\balance

\end{document}